\section{Angular Momentum Operators}
\subsection{Generators}\label{generators}

So far, we've considered several Hermitian operators: the position operator $\hat{x}$, the momentum operator $\hat{p}$, and the energy operator, or Hamiltonian, $\hat{H}$. As we begin our work in $3$ dimensions, we now discuss yet another operator: angular momentum. Like momentum, it is a vector operator. It will have three components, each of which is Hermitian and thus represents a measurable quantity. Although the definition of the angular momentum operator will be motivated by its classical counterpart, we will see unexpected quantum properties.

It turns out that the momentum operator has something to do with translations; specifically, consider the operator $$e^{\frac{i\hat{p}a}{\hbar}},$$where $a$ is a constant with units of length, acting on a wavefunction $\psi$: $$e^{\frac{i\hat{p}a}{\hbar}}\psi(x)=e^{a\frac{d}{dx}}\psi(x)=\sum_{n=0}^\infty\frac{1}{n!}\bigg(a\frac{d}{dx}\bigg)^n\psi(x).$$Expanding the sum gives $\psi(x)+a\frac{d\psi}{dx}+\frac{a^2}{2!}\frac{d^2\psi}{dx^2}+\frac{a^3}{3!}\frac{d^3\psi}{dx^3}+\ldots$ Note that this is simply the Taylor series expansion of $\psi(x+a)$. This means that the operator $e^{\frac{i\hat{p}a}{\hbar}}$ moves the wavefunction a displacement of $-a$. We say that the momentum operator \textbf{generates} translations. Similarly, we will show that the \textit{angular} momentum operator generates \textit{rotations}.

\begin{framedrmk*}[Types of angular momentum]
    There are two types of angular momentum: \textbf{orbital} and \textbf{spin}. Orbital angular momentum is the familiar kind that occurs when a particle rotates around some fixed point. But there is also spin angular momentum, which is carried by point particles that somehow possess angular momentum despite having no apparent size. Much of the mathematics of angular momentum is valid for both orbital and spin angular momentum.
\end{framedrmk*}

\subsection{Schrödinger Equation for a Central Potential}\label{schrodinger_central_potential}

Let's begin our analysis of angular momentum by recalling that in three dimensions the usual position and momentum operators are vector operators: \begin{align*}\hat{\mathbf{p}}&= (\hat{p}_x,\hat{p}_y,\hat{p}_z)=\frac{\hbar}{i}\nabla=\frac{\hbar}{i}\bigg(\frac{\partial}{\partial x},\frac{\partial}{\partial y},\frac{\partial}{\partial z}\bigg),\\\hat{\mathbf{x}}&=(\hat{x},\hat{y},\hat{z}).\end{align*}The commutation relations are $$[\hat{x},\hat{p}_x]= [\hat{y},\hat{p}_y]=[\hat{z},\hat{p}_z]=i\hbar.$$All other commutators involving the position and momentum operators are zero.

Consider a particle represented by a three-dimensional wavefunction $\psi(x,y,z)$ moving in a potential $V(\mathbf{r})$. The three-dimensional Schrödinger equation takes the form $$-\frac{\hbar^2}{2m}\nabla^2\psi(\mathbf{r})+V(\mathbf{r})\psi(\mathbf{r})=E\psi(\mathbf{r}).$$We have a \textbf{central potential} if $V$ is spherically symmetric -- that is, if $V(\mathbf{r})=V(r)$. Then the equipotential surfaces are spheres centered at the origin. For a central potential, the above equation is $$-\frac{\hbar^2}{2m}\nabla^2\psi(\mathbf{r})+V(r)\psi(\mathbf{r})=E\psi(\mathbf{r}).$$Note that the wavefunction is *not* necessarily spherically symmetric; in fact, it will only be for the simplest kinds of solutions. Given the rotational symmetry of the potential, however, we express the following equations in spherical coordinates. The Laplacian is given by $$\nabla^2\psi=(\nabla\cdot\nabla)\psi=\frac{1}{r}\frac{\partial^2}{\partial r^2}(r\psi)+\frac{1}{r^2}\bigg(\frac{1}{\sin\theta}\frac{\partial}{\partial\theta}\sin\theta\frac{\partial}{\partial\theta}+\frac{1}{\sin^2\theta}\frac{\partial^2}{\partial\phi^2}\bigg)\psi.$$Therefore the Schrödinger equation for a particle in a central potential becomes $$-\frac{\hbar^2}{2m}\bigg[\frac{1}{r}\frac{\partial^2}{\partial r^2}(r\psi)+\frac{1}{r^2}\bigg(\frac{1}{\sin\theta}\frac{\partial}{\partial\theta}\sin\theta\frac{\partial}{\partial\theta}+\frac{1}{\sin^2\theta}\frac{\partial^2}{\partial\phi^2}\bigg)\bigg]\psi+V(r)\psi=E\psi.$$In the work that follows, we will seek to establish the following two statements.
\begin{framedprop}\label{12.1}
    The angular dependent piece of the Laplacian operator $\nabla^2$ is proportional to the magnitude squared of the angular momentum operator: $$\frac{1}{\sin\theta}\frac{\partial}{\partial\theta}\bigg(\sin\theta\frac{\partial}{\partial\theta}\bigg)+\frac{1}{\sin^2\theta}\frac{\partial^2}{\partial\phi^2}=-\frac{\hat{\mathbf{L}}^2}{\hbar^2},$$where $$\hat{\mathbf{L}}^2=\hat{L}_x\hat{L}_x+\hat{L}_y\hat{L}_y+\hat{L}_z\hat{L}_z.$$This will imply that the Schrödinger equation becomes $$-\frac{\hbar^2}{2m}\frac{1}{r}\frac{\partial}{\partial r^2}(r\psi)+\frac{\mathbf{L}^2}{2mr^2}+V(r)\psi=E\psi.$$
\end{framedprop}
\begin{framedprop}\label{12.2}
    The Schrödinger equation for a central potential is valid for the \textit{two-body} problem when the potential $V(\mathbf{r}_1,\mathbf{r}_2)$ is only a function of the distance between the objects -- that is, when $$V(\mathbf{r}_1,\mathbf{r}_2)=V(|r_1-r_2|).$$This is true for the electrostatic potential energy between the proton and the electron of a hydrogen atom. Thus, we will be able to treat the hydrogen atom as a central potential problem.
\end{framedprop}

\subsection{Angular Momentum Algebra}\label{angular_momentum_algebra}

Classically, we are familiar with the definition of the angular momentum as the cross product of $\mathbf{r}$ and $\mathbf{p}$. We have \begin{align*}\mathbf{L}&=(L_x,L_y,L_z)=\mathbf{r}\times\mathbf{p},\\ L_x&=yp_z-zp_y,\\ L_y&=zp_x-xp_z, \\ L_z &= xp_y-yp_x. \end{align*}We use these relations to define the quantum angular momentum operator $\hat{\mathbf{L}}$ and its components: \begin{align*}\hat{\mathbf{L}}&=(\hat{L}_x,\hat{L}_y,\hat{L}_z),\\ \hat{L}_x&=\hat{y}\hat{p}_z-\hat{z}\hat{p}_y,\\ \hat{L}_y&=\hat{z}\hat{p}_x-\hat{x}\hat{p}_z, \\ \hat{L}_z &= \hat{x}\hat{p}_y-\hat{y}\hat{p}_x. \end{align*}Note that the order of the terms in each product doesn't matter, since a coordinate and a momentum along different axes commute. It follows that these operators are Hermitian. For example, recalling that $(AB)^\dagger=B^\dagger A^\dagger$, we have $$(\hat{L}_x)^\dagger=(\hat{y}\hat{p}_z-\hat{z}\hat{p}_y)^\dagger=(\hat{y}\hat{p}_z)^\dagger-(\hat{z}\hat{p}_y)^\dagger=\hat{p}_z^\dagger\hat{y}^\dagger-\hat{p}_y^\dagger\hat{z}^\dagger.$$Since coordinates and momenta are Hermitian operators, we have $$(\hat{L}_x)^\dagger=\hat{p}_z\hat{y}-\hat{p}_y\hat{z}=\hat{y}\hat{p}_z-\hat{z}\hat{p}_y=\hat{L}_x.$$Therefore, all the angular momentum operators represent are Hermitian and observables, as they should.

Given a set of Hermitian operators, a natural question is that of their commutators. This computation enables us to see if we can measure them simultaneously. Let us compute the commutator of $\hat{L}_x$ with $\hat{L}_y$: $$[\hat{L}_x,\hat{L}_y]=[\hat{y}\hat{p}_z-\hat{z}\hat{p}_y,\hat{z}\hat{p}_x-\hat{x}\hat{p}_z].$$We now see that only the $\hat{z}$ and $\hat{p}_z$ terms fail to commute. Therefore, \begin{align*}[\hat{L}_x,\hat{L}_y]&=[\hat{y}\hat{p}_z,\hat{z}\hat{p}_x]+[\hat{z}\hat{p}_y,\hat{x}\hat{p}_z]\\ &=[\hat{y}\hat{p}_z,\hat{z}]\hat{p}_x+\hat{x}[\hat{z}\hat{p}_y,\hat{p}_z] \\ &=\hat{y}[\hat{p}_z,\hat{z}]\hat{p}_x+\hat{x}[\hat{z},\hat{p_z}]\hat{p}_y \\ &=\hat{y}(-i\hbar)\hat{p}_x+\hat{x}(i\hbar)\hat{p}_y \\ &=i\hbar(\hat{x}\hat{p}_y-\hat{y}\hat{p}_x),\end{align*}which we recognize as $i\hbar\hat{L}_z$! Therefore, $$[\hat{L}_x,\hat{L}_y]=i\hbar\hat{L}_z.$$The basic commutation relations are completely cyclic, so we have \begin{align*}[\hat{L}_x,\hat{L}_y]&=i\hbar\hat{L}_z,\\ [\hat{L}_y,\hat{L}_z] &= i\hbar\hat{L}_x, \\ [\hat{L}_z,\hat{L}_x] &= i\hbar\hat{L}_y.\end{align*}
The set of angular momentum operators obeying these commutation relations is known as the \textbf{quantum algebra of angular momentum}. 

\begin{framedrmk*}[Lie algebras]
    The angular momentum algebra is so important that it appears in all fields of physics and mathematics. It is \href{https://www.sjsu.edu/faculty/watkins/SU2.htm}{related to the algebra of generators of the group $SU(2)$}, the orthogonal group in 3 dimensions, and quaternions. In fact, the subject of Lie algebras seeks to classify all possible consistent commutation relations, and this is the first non-trivial example.
\end{framedrmk*}

The $\hat{L}$ operators are sometimes referred to as \textit{orbital} angular momentum operators, to distinguish them from \textit{spin} angular momentum operators $\hat{S}_x$, $\hat{S}_y$, and $\hat{S}_z$. These latter operators cannot be written in terms of coordinates and momenta. However, they still satisfy exactly the same algebra as the $\hat{L}$ operators: \begin{align*}[\hat{S}_x,\hat{S}_y]&=i\hbar\hat{S}_z,\\ [\hat{S}_y,\hat{S}_z] &= i\hbar\hat{S}_x, \\ [\hat{S}_z,\hat{S}_x] &= i\hbar\hat{S}_y.\end{align*}

\subsection{Commuting Observables}\label{commuting_observables}

We have now defined the angular momentum operators $\hat{L}_x$, $\hat{L}_y$, and $\hat{L}_z$, and seen that they are all Hermitian. We now ask the following question.

\begin{framedquest*}
    Can we have simultaneous eigenstates of $\hat{L}_x,\,\hat{L}_y$, and $\hat{L}_z?$
\end{framedquest*}

It turns out that similarly to the way we cannot have a simultaneous eigenstate of $\hat{x}$ and $\hat{p}$, we also cannot have a simultaneous eigenstate of two distinct angular momentum operators. We demonstrate this as follows.

Suppose there exists a wavefunction $\phi_0$ which is simultaneously an eigenstate of $\hat{L}_x$ and $\hat{L}_y$: \begin{align*}\hat{L}_x\phi_0&=\lambda_x\phi_0,\\ \hat{L}_y\phi_0&=\lambda_y\phi_0.\end{align*}However, by the commutator relations, we have $$i\hbar\hat{L}_z\phi_0=[\hat{L}_x,\hat{L}_y]\phi_0=\hat{L}_x\hat{L}_y\phi_0-\hat{L}_y\hat{L}_x\phi_0=(\lambda_x\lambda_y-\lambda_y\lambda_x)\phi_0=0.$$Therefore, $\hat{L}_z\phi_0=0$. However, we also must have that \begin{align*}[\hat{L}_y,\hat{L}_z]\phi_0&=0=i\hbar\hat{L}_x\phi_0=i\hbar\lambda_x\phi_0\implies\lambda_x=0,\\ [\hat{L}_z,\hat{L}_x]\phi_0&=0=i\hbar\hat{L}_y\phi_0=i\hbar\lambda_y\phi_0\implies\lambda_y=0.\end{align*}Therefore, if $\phi_0$ is a simultaneous eigenstate of two distinct angular momentum operators, it is annihilated by all of them. It is impossible to find states that are nontrivial simultaneous eigenstates of two angular momentum operators.

We still wish to find some operator that commutes with $\hat{L}_x$. Since $\hat{L}_x$ represents some rotation in 3 dimensions, we can intuitively expect any operator that commutes with it to be rotationally symmetric. It turns out that $\hat{\mathbf{L}}^2$, as defined in Proposition \ref{12.1}, works! We can quickly check this: \begin{align*}[\hat{L}_x,\hat{\mathbf{L}}^2]&=[\hat{L}_x,\hat{L}_x\hat{L}_x+\hat{L}_y\hat{L}_y+\hat{L}_z\hat{L}_z] \\ &=[\hat{L}_x,\hat{L}_y\hat{L}_y+\hat{L}_z\hat{L}_z]\\ &=[\hat{L}_x,\hat{L}_y]\hat{L}_y+\hat{L}_y[\hat{L}_x,\hat{L}_y]+[\hat{L}_x,\hat{L}_z]\hat{L}_z+\hat{L}_z[\hat{L}_x,\hat{L}_z]\\ &=i\hbar\hat{L}_z\hat{L}_y+i\hbar\hat{L}_y\hat{L}_z-i\hbar\hat{L}_y\hat{L}_z-i\hbar\hat{L}_z\hat{L}_y\\ &=0.\end{align*}Thus, we can find simultaneous eigenstates of both $\hat{L}_x$ and $\hat{\mathbf{L}}^2$. And there's nothing special about $\hat{L}_x$ --- the operator $\hat{\mathbf{L}}^2$ commutes with all three angular momentum operators. Such an operator is called a \textbf{Casimir operator}.

\begin{framedrmk*}[Simultaneous eigenstates]
    It is a general theorem of linear algebra that given any two operators that commute, it is possible to find simultaneous eigenstates of both operators.
\end{framedrmk*}

\begin{framedrmk*}[Why do we care?]
    You might wonder why we even care about simultaneous eigenstates. In any physical context, we want to know the most we can about the system we are given. A state that is a simultaneous eigenstate of energy and momentum, for example, is more helpful for us to study, since we will know two different quantities about our system. In general, physicists search for the maximal set of commuting observables to study.
\end{framedrmk*}

\subsection{Spherical Coordinates}\label{spherical_coordinates}

To better understand the angular momentum operators, we write them in spherical coordinates. For this, we need the relations between $(r,\theta,\phi)$ and the Cartesian coordinates $(x,y,z)$: \begin{align*}x&= r\sin\theta\cos\phi  &r= \sqrt{x^2+y^2+z^2}\\ y&=r\sin\theta\sin\phi &\theta=\cos^{-1}(\frac{z}{r}) \\ z&=r\cos\theta &\phi=\tan^{-1}(\frac{y}{x})\end{align*}We mentioned in Section \ref{generators} that angular momentum operators generate rotations. In spherical coordinates, rotations about the $z$ axis are the simplest: they change $\phi$ but leave $\theta$ invariant. (Rotations about the $x$ and $y$ axes change both $\theta$ and $\phi$.) We therefore express $\hat{L}_z$ in spherical coordinates: $$\hat{L}_z=\hat{x}\hat{p}_y-\hat{y}\hat{p}_x=\frac{\hbar}{i}\bigg(x\frac{\partial}{\partial y}-y\frac{\partial}{\partial x}\bigg).$$Also, note that applying the chain rule to $\frac{\partial}{\partial \phi}$ gives $$\frac{\partial}{\partial\phi}=\frac{\partial x}{\partial\phi}\frac{\partial}{\partial x}+\frac{\partial y}{\partial\phi}\frac{\partial}{\partial y}+\frac{\partial z}{\partial\phi}\frac{\partial}{\partial z}=x\frac{\partial}{\partial y}-y\frac{\partial}{\partial x}$$using the relations above. We conclude that $$\boxed{\hat{L}_z=\frac{\hbar}{i}\frac{\partial}{\partial\phi}.}$$The similarity between this expression and $\hat{p}=\frac{\hbar}{i}\frac{\partial}{\partial x}$ confirms the interpretation that $\hat{L}_z$ generates changes in $\phi$, which are rotations about the $z$-axis. A lengthier calculation for the other angular momentum operators proves Proposition \ref{12.1}: $$-\frac{\hat{\mathbf{L}}^2}{\hbar^2}=\frac{1}{\sin\theta}\frac{\partial}{\partial\theta}\bigg(\sin\theta\frac{\partial}{\partial\theta}\bigg)+\frac{1}{\sin^2\theta}\frac{\partial^2}{\partial\phi^2}.$$

\section{Eigenstates of Angular Momentum}
\subsection{Quantization of Angular Momentum}\label{quantization_angular_momentum}

We demonstrated above that $\hat{\mathbf{L}}^2$ commutes with all three angular momentum operators. Considering the simplicity of $\hat{L}_z$, we now aim to construct the simultaneous eigenstates of these two operators. They will be functions of $\theta$ and $\phi$, and we will call them $\psi_{lm}(\theta,\phi)$. The conditions for $\psi_{lm}$ to be an eigenstate of both operators are \begin{align*}\hat{L}_z\psi_{lm}&=\hbar m\,\psi_{lm},&m\in\mathbb{R}\\ \hat{\mathbf{L}}^2\psi_{lm}&=\hbar^2l(l+1)\psi_{lm},& l\in\mathbb{R}.\end{align*}

Both operators are Hermitian, so the eigenvalues are real. We also add $\hbar$ factors to make $m$ and $l$ dimensionless. Furthermore, the eigenvalues of $\hat{\mathbf{L}}^2$ can't be negative. For this, we first claim that $$(\psi,\hat{\mathbf{L}}^2\psi)\geq0$$for all $\psi$, which we can prove by simply expanding: $$(\psi,\hat{\mathbf{L}}^2)=(\psi,\hat{L}_x^2\psi)+(\psi,\hat{L}_y^2\psi)+(\psi,\hat{L}_z^2\psi).$$By Hermiticity, $$(\psi,\hat{\mathbf{L}}^2)=(\hat{L}_x\psi,\hat{L}_x\psi)+(\hat{L}_y\psi,\hat{L}_y\psi)+(\hat{L}_z\psi,\hat{L}_z\psi)\geq 0,$$since each of the three terms is non-negative. Now taking $\psi$ to be a normalized eigenfunction with eigenvalue $\lambda$ gives $(\psi,\lambda\psi)=\lambda\geq 0$, as desired. This allows us to express $\lambda=l(l+1)\geq-\frac{1}{4}$ for $l\in\mathbb{R}$.

Let's now solve the first eigenvalue equation using the spherical representation of the $\hat{L}_z$ operator: $$\frac{\hbar}{i}\frac{\partial}{\partial \phi}\psi_{lm}=\hbar m\,\psi_{lm}\rightarrow \frac{\partial}{\partial \phi}\psi_{lm}=im\,\psi_{lm}.$$This determines the $\phi$ dependence of the solution, and we write $$\boxed{\psi_{lm}(\theta,\phi)=e^{im\phi}P_l^m(\theta),}$$where the function $P_l^m(\theta)$ captures the still undetermined $\theta$ dependence of the eigenfunction $\psi_{lm}$. We will also require that $\psi_{lm}$ is periodic in $\phi$: $$\psi_{lm}(\theta,\phi+2\pi)=\psi_{lm}(\theta,\phi),$$which requires that $$e^{im(\phi+2\pi)}=e^{im\phi}\rightarrow e^{2\pi im}=1.$$Thus, $m$ is an integer, and we've reached our first quantization: $$\boxed{m\in\mathbb{Z}.}$$The second eigenvalue equation is more complicated. Using the expression for $\hat{\mathbf{L}}^2$, we have $$-\hbar^2\bigg(\frac{1}{\sin\theta}\frac{\partial}{\partial\theta}\bigg(\sin\theta\frac{\partial}{\partial\theta}\bigg)+\frac{1}{\sin^2\theta}\frac{\partial^2}{\partial\phi^2}\bigg)\psi_{lm}=\hbar^2l(l+1)\psi_{lm}.$$Multiplying through by $\sin^2\theta$ and cancelling the $\hbar^2$ factors, $$\bigg(\sin\theta\frac{\partial}{\partial\theta}\bigg(\sin\theta\frac{\partial}{\partial\theta}\bigg)+\frac{\partial^2}{\partial\phi^2}\bigg)\psi_{lm}=-l(l+1)\sin^2\theta\psi_{lm}.$$Using $\psi_{lm}(\theta,\phi)=e^{im\phi}P_l^m(\theta)$, $$\sin\theta\frac{d}{d\theta}\bigg(\sin\theta\frac{dP_l^m}{d\theta}\bigg)-m^2P_l^m=-l(l+1)P_l^m\sin^2\theta.$$Equivalently, $$\sin\theta\frac{d}{d\theta}\bigg(\sin\theta\frac{dP_l^m}{d\theta}\bigg)+[l(l+1)\sin^2\theta-m^2]P_l^m=0.$$Since $\theta$ ranges from $0$ to $\pi$ in spherical coordinates, we parameterize it using $x\equiv\cos\theta$. This gives $$\frac{d}{d\theta}=\frac{dx}{d\theta}\frac{d}{dx}=-\sin\theta\frac{d}{dx}\rightarrow \sin\theta\frac{d}{dx}=-(1-x^2)\frac{d}{dx}.$$With this substitution, the differential equation becomes $$(1-x^2)\frac{d}{dx}\bigg[(1-x^2)\frac{dP_l^m}{dx}\bigg]+[l(l+1)(1-x^2)-m^2]P_l^m(x)=0.$$Dividing by $(1-x^2)$, we finally arrive at the final form: $$\boxed{\frac{d}{dx}\bigg[(1-x^2)\frac{dP_l^m}{dx}\bigg]+\bigg[l(l+1)-\frac{m^2}{1-x^2}\bigg]P_l^m(x)=0.}$$The $P_l^m(x)$ given by this equation are called the \textbf{associated Legendre functions}. We know almost nothing about them at this point; only that $m$ is an integer. To find out about $l$, we begin by considering the case when $m=0$. For brevity, we write $P_l(x)\equiv P_l^0(x)$. We must have $$\frac{d}{dx}\bigg[(1-x^2)\frac{dP_l}{dx}\bigg]+l(l+1)P_l(x)=0.$$This is the \textbf{Legendre differential equation}. We find a series solution by writing $$P_l(x)=\sum_{k=0}^\infty a_kx^k.$$Substituting this into the differential equation, we find that the vanishing of the coefficient of $x^k$ requires $$(k+1)(k+2)a_{k+2}+[(l(l+1)-k(k+1)]a_k=0.$$Equivalently, we have $$\frac{a_{k+2}}{a_k}=-\frac{l(l+1)-k(k+1)}{(k+1)(k+2)}.$$As $k\to\infty$, $P_l$ diverges at $x=\pm1$. Therefore, the series must terminate, and so $l(l+1)$ must equal $k(k+1)$ for some integer $k\geq 0$. We can simply pick $l=k$ so that $a_{k+2}=0$, making $P_k(x)$ a degree-$k$ polynomial. Thus, $l$ can be any natural number, which gives us our second quantization: $$\boxed{l\in\mathbb{N}_0.}$$

\begin{frameddefn}[Legendre polynomials]\label{legendre_poly}
    The \textbf{Legendre polynomials} $P_l(x)$ are given by the \textbf{Rodriguez formula}: $$P_l(x)=\frac{1}{2^ll!}\big(\frac{d}{dx}\big)^l(x^2-1)^l,$$and have a nice generating function $$\sum_{l=j}^\infty P_l(x)s^l=\frac{1}{\sqrt{1-2xs+s^2}}.$$The first few Legendre polynomials are \begin{align*}P_0(x)&=1,\\ P_1(x) &=x,\\ P_2(x) &= \frac{1}{2}(3x^2-1).\end{align*}
\end{frameddefn}

\subsection{Spherical Harmonics}\label{spherical_harmonics}

Let's rewrite the differential equation for the associated Legendre functions: $$\frac{d}{dx}\bigg[(1-x^2)\frac{dP_l^m}{dx}\bigg]+\bigg[l(l+1)-\frac{m^2}{1-x^2}\bigg]P_l^m(x)=0.$$Notice that this equation depends only on $m^2$, so we can define the solutions for $m$ and $-m$ to be the same. In fact, it turns out that the solutions to this equation are $$P_l^m(x)=(1-x^2)^{\frac{|m|}{x}}\bigg(\frac{d}{dx}\bigg)^{|m|}P_l(x)$$Recall that the Legendre polynomials $P_l(x)$ are of the form $$P_l(x)=\frac{1}{2^ll!}\big(\frac{d}{dx}\big)^l(x^2-1)^l,$$which is a polynomial of degree $l$. Therefore, the above solutions are non-zero only for $|m|\leq l$. We thus have solutions for $$\boxed{-l\leq m\leq l}$$It is possible to prove that there are no solutions to the above differential equation other than the $P_l^m(x)$ given above. 

Recall that $P_l^m(\theta)$ gives the angular dependence of the angular momentum eigenstates. One can think of the $\psi_{lm}$ eigenfunctions as first determined by the whole number $l$, and for a fixed $l$, the integer $m$.

What \textit{are} the $\psi_{lm}$ eigenfunctions? With suitable normalization, they are called the \textbf{spherical harmonics} $Y_{l,m}(\theta,\phi)$. For $m\geq 0$, we define them as follows.

\begin{frameddefn}[Spherical harmonics]
    For $m\geq 0$ and $m\leq l$, we define the \textbf{spherical harmonic} $$Y_{l,m}(\theta,\phi)=\sqrt{\frac{2l+1}{4\pi}\frac{(l-m)!}{(l+m)!}}(-1)^me^{im\phi]}P_l^m(\cos\theta).$$For $m<0$ and $m>-l$, we take $$Y_{l,m}=(-1)^m[Y_{l,-m}(\theta,\phi)]^*.$$The first few spherical harmonics are \begin{align*}Y_{0,0}(\theta,\phi)&=\frac{1}{\sqrt{4\pi}},\\ Y_{1,\pm1}(\theta,\phi)&=\mp\sqrt{\frac{3}{8\pi}}e^{\pm i\phi}\sin\theta,\\ Y_{1,0}(\theta,\phi)&=\sqrt{\frac{3}{4\pi}}\cos\theta.\end{align*}
\end{frameddefn}

Since they are eigenstates of Hermitian operators with different eigenvalues, spherical harmonics with different subscripts are automatically orthogonal. Furthermore, they form an orthonormal set with respect to integration over solid angle.

\begin{framedrmk*}[Integration over solid angle]
    Integrating over solid angle $\Omega$ is like integrating over the surface of a sphere of radius 1. It can be written in many forms: $$\int d\Omega=\int_0^{2\pi}d\phi\int_0^\pi d\theta\sin\theta=\int_0^{2\pi}d\phi\int_{-1}^1d(\cos\theta).$$
\end{framedrmk*}

The statement that the spherical harmonics form an orthonormal set with respect to this integration means that $$\int d\Omega\, Y^*_{l',m'}(\theta,\phi)Y_{l,m}(\theta,\phi)=\delta_{l,l'}\delta_{m,m'}.$$In other words, all the spherical harmonics are normalized, and the inner product of two different spherical harmonics is $0$.

\subsection{The Radial Equation}\label{radial_equation}

We now return to the Schrödinger equation $$-\frac{\hbar^2}{2m}\nabla^2\psi+V(r)\psi=E\psi.$$By the definition of the Laplacian, we have $$-\frac{\hbar^2}{2m}\bigg[\frac{1}{r}\frac{d^2}{dr^2}r\psi-\frac{1}{\hbar^2r^2}L^2\psi\bigg]+V(r)\psi=E\psi$$We want to consider the radial component of our solutions, so we try factorizing $\psi$. Specifically, we let $$\psi(\mathbf{r})=R_{E}(r)Y_{l,m}(\theta,\phi).$$Recall that  $$\hat{\mathbf{L}}^2Y_{lm}=\hbar^2l(l+1)Y_{lm}. $$Substituting, we have $$-\frac{\hbar^2}{2m}\bigg[\frac{1}{r}\frac{d^2}{dr^2}(rR_E)-\frac{l(l+1)}{r^2}R_E\bigg]+V(r)R_E=ER_E$$Multiplying both sides by $r$, $$-\frac{\hbar^2}{2m}\bigg[\frac{d^2}{dr^2}(rR_E)-\frac{l(l+1)}{r^2}(rR_E)\bigg]+V(r)(rR_E)=E(rR_E)$$Let $u_E(r)=rR_E(r)$. Then the above equation becomes $$\boxed{-\frac{\hbar^2}{2m}\frac{d^2u_E}{dr^2}+\bigg[\frac{\hbar^2l(l+1)}{2mr^2}+V(r)\bigg]u_E=Eu_E.}$$This is called the \textbf{radial equation}, and expresses the radial component of the Schrödinger equation. It is such an important formulation of the Schrödinger equation that I will rewrite it in a box.

\begin{framedthm}[Radial equation]
    $$-\frac{\hbar^2}{2m}\frac{d^2u_E}{dr^2}+\bigg[\frac{\hbar^2l(l+1)}{2mr^2}+V(r)\bigg]u_E=Eu_E.$$
\end{framedthm}

\subsection{Effective Potential}\label{effective_potential}

The radial equation looks like the familiar time-independent Schrödinger equation in one dimension, but with an \textit{effective potential} $$V_{eff}(r)=V(r)+\frac{\hbar^2l(l+1)}{2mr^2}.$$In the case when $l=0$, that is when the particle has no angular momentum, the effective potential is simply $V(r)$. Note that $V_{eff}(r)$ is independent of $m$. Intuitively, $l$ encodes the total angular momentum of a particle, while $m$ encodes the $z$-component of the angular momentum.

Recall our decomposition of the wavefunction $$\psi(r,\theta,\phi)=\frac{u_{E,l}(r)}{r}Y_{l,m}(\theta,\phi).$$We have added a subscript $l$ to $u_E(r)$, since $u_E(r)$ depends on the effective potential, which depends on $l$.

The normalization condition for $\psi$ gives $$\int|\psi|^2d^3x=\int_0^\infty r^2dr\int d\Omega\, \frac{|u|^2}{r^2}Y_{lm}^*Y_{lm}=\int_0^\infty dr\,|u(r)|^2=1.$$We can see that the three-dimensional wavefunctions $\psi$ are normalized if and only if the one-dimensional wavefunctions $u$ are. Notice that only $r>0$ is allowed, so we should also consider the possible behavior of the wavefunction when $r$ approaches $0$.

We can learn about this behavior under the reasonable assumption that the centrifugal barrier dominates the potential as $r\to 0$. In this case, the centrifugal term must be roughly canceled out by the differential term: $$-\frac{\hbar^2}{2m}\frac{d^2u}{dr^2}+\frac{\hbar^2l(l+1)}{2mr^2}u\approx 0.$$Simplifying, $$\frac{d^2u}{dr^2}\approx\frac{l(l+1)}{r^2}u.$$The solutions are of the form $u=r^{l+1}$ and $u=r^{-l}$. Since we are considering the regime as $r\to0$, the second expression diverges and cannot be normalized. Therefore, we have that $$u_{E,l}(r)\sim r^{l+1} \text{ as }r\to0.$$Regardless of the value of $l$, $u_{E,l}$ vanishes at $r=0$. Effectively, there is an infinite wall at $r=0$, which represents the impossibility of extending $r$ to negative values. On the other hand, the radial dependence of the wavefunction $\psi$ then goes like $\frac{r^{l+1}}{r}=r^l$ for $r\approx 0$. This allows for a constant non-zero wavefunction at $r=0$ if $l=0$. For $l\neq 0$, an angular momentum barrier prevents the particle from reaching the origin.